% Демонстрация 7: «Линейная зависимость векторов на плоскости» (Лекции 15-16)
\section*{Демонстрация 7. «Линейная зависимость векторов на плоскости»\\ \normalsize\textmd{Файл \texttt{demos/07-linear-dependence.html} \quad(к Лекциям~15--16)}}

\subsection*{Какие факты лекций иллюстрирует}

Демонстрация связывает воедино два смежных сюжета: \textbf{линейную
(не)зависимость} пары векторов (Лекция~15) и \textbf{базис плоскости с
координатами вектора} (Лекция~16). Студент перетаскиванием задаёт два
базисных вектора $\bar{e}_1$, $\bar{e}_2$ и целевой вектор $\bar{x}$;
программа решает систему, выводит координаты разложения $\alpha_1$,
$\alpha_2$, рисует параллелограмм разложения, а при коллинеарности базиса
подсвечивает определитель $\Delta = 0$ и сообщение «базис вырожден».

\begin{itemize}
  \item \textbf{Критерий линейной (не)зависимости (Лекция~15).} Базис
    плоскости линейно независим, если нулю равна только тривиальная
    комбинация:
    \[
      \lambda_1 \bar{e}_1 + \lambda_2 \bar{e}_2 = \bar 0
      \iff \lambda_1 = \lambda_2 = 0.
    \]
    Линейная зависимость, наоборот, означает существование
    нетривиального набора: $\exists\,(\lambda_i) \neq 0:\
    \sum_i \lambda_i \bar v_i = \bar 0$. Сдвигая $\bar{e}_2$ к одной
    прямой с $\bar{e}_1$, студент воспроизводит ровно тот случай, когда
    один вектор выражается через другой ($\bar{e}_2 = \lambda\bar{e}_1$),
    --- это \textbf{критерий коллинеарности} из Лекции~15: два вектора
    линейно зависимы тогда и только тогда, когда они коллинеарны.

  \item \textbf{Базис и размерность плоскости (Лекция~16).} Любые два
    неколлинеарных вектора на плоскости линейно независимы и образуют
    базис, а размерность плоскости равна двум: третий вектор $\bar{x}$
    заведомо выражается через $\bar{e}_1$, $\bar{e}_2$. Демонстрация
    делает это утверждение осязаемым --- $\bar{x}$ всегда удаётся
    разложить, пока базис невырожден.

  \item \textbf{Единственность разложения (Лекция~16).} Пока $\Delta\neq
    0$, координаты $\alpha_1$, $\alpha_2$ восстанавливаются однозначно ---
    это иллюстрация теоремы о разложении вектора по базису и критерия
    единственности разложения из Лекции~15. Как только базис вырождается,
    разложение либо невозможно (если $\bar{x}$ вне общей прямой), либо
    неоднозначно (если $\bar{x}$ на ней), и панель координат гаснет.

  \item \textbf{Определитель как критерий базиса.} Базис плоскости
    невырожден тогда и только тогда, когда
    \[
      \Delta =
      \begin{vmatrix}
        e_{1x} & e_{2x} \\
        e_{1y} & e_{2y}
      \end{vmatrix} \neq 0.
    \]
    Это двумерный аналог проверки тройки векторов в $V_3$ через
    смешанное произведение (определитель $3\times 3$) из Лекции~15 и
    условие невырожденности матрицы перехода из Лекции~16.

  \item \textbf{Формулы Крамера для координат.} Координаты разложения
    программа считает по правилу Крамера:
    \[
      \alpha_1 = \frac{1}{\Delta}
      \begin{vmatrix}
        x_x & e_{2x} \\
        x_y & e_{2y}
      \end{vmatrix},
      \qquad
      \alpha_2 = \frac{1}{\Delta}
      \begin{vmatrix}
        e_{1x} & x_x \\
        e_{1y} & x_y
      \end{vmatrix}.
    \]
    Знаменатель $\Delta$ наглядно объясняет, почему при вырождении базиса
    координаты не определены: на нуль делить нельзя.
\end{itemize}

\begin{remarkIT}{Где это работает за пределами лекции}{demo07_apps}
  Разложение вектора по базису плоскости --- упрощённая модель сразу
  нескольких прикладных сюжетов:
  \begin{itemize}
    \item \textbf{Системы координат в графике.} Пара $\bar{e}_1$,
      $\bar{e}_2$ --- это локальный базис (UV-координаты текстуры,
      экранные оси спрайта), а $\alpha_1$, $\alpha_2$ --- координаты
      точки в нём. Вырождение базиса соответствует «схлопыванию» системы
      координат в линию.
    \item \textbf{Смена базиса и матрица перехода.} Столбцы из координат
      $\bar{e}_1$, $\bar{e}_2$ --- это матрица перехода $C$ (Лекция~16);
      решение системы --- умножение $\bar{x}$ на $C^{-1}$. Условие
      $\Delta\neq 0$ --- та самая невырожденность $C$, гарантирующая
      обратимость преобразования координат.
    \item \textbf{Признаки в ML и мультиколлинеарность.} Коллинеарные
      $\bar{e}_1$, $\bar{e}_2$ --- модель мультиколлинеарности столбцов
      матрицы признаков: $\Delta\to 0$ означает почти вырожденную
      $X^TX$ и неустойчивую регрессию.
    \item \textbf{Ранг матрицы.} Невырожденный базис --- это матрица
      $2\times 2$ полного ранга~$2$; коллинеарность роняет ранг до~$1$.
    \item \textbf{Разреженные базисы.} Удачный базис делает координаты
      $\alpha_1$, $\alpha_2$ простыми (многие нулевые) --- идея,
      лежащая в основе сжатия сигналов и разреженного кодирования.
  \end{itemize}
\end{remarkIT}

\subsection*{Примерный сценарий использования на занятии}

\begin{enumerate}
  \item \textbf{Когда включить.} Демонстрацию запускают на стыке двух
    лекций: после критерия коллинеарности (Лекция~15) и определения
    базиса и координат (Лекция~16), как только на доске выписано
    разложение $\bar{x} = \alpha_1\bar{e}_1 + \alpha_2\bar{e}_2$.

  \item \textbf{Невырожденный базис.} Оставьте стартовую конфигурацию
    ($\bar{e}_1$, $\bar{e}_2$ неколлинеарны) и подвигайте только
    $\bar{x}$. Студенты видят, как координаты $\alpha_1$, $\alpha_2$ и
    параллелограмм разложения отслеживают положение $\bar{x}$. Вопрос:
    \emph{«Единственно ли разложение?»} --- да, пока базис фиксирован,
    каждой точке отвечает ровно одна пара координат (теорема о
    разложении по базису).

  \item \textbf{Геометрия координат.} Попросите предсказать координаты
    до отпускания мыши: если $\bar{x}$ совпал с концом $\bar{e}_1$, то
    $\alpha_1 = 1$, $\alpha_2 = 0$. Это закрепляет смысл координат как
    «доли» базисных векторов в разложении.

  \item \textbf{Движение к вырождению.} Медленно ведите $\bar{e}_2$ к
    прямой, на которой лежит $\bar{e}_1$, и следите за панелью:
    определитель $\Delta$ плавно убывает к нулю, а параллелограмм
    сплющивается в отрезок. Вопрос: \emph{«Что значит $\Delta = 0$
    геометрически?»} --- площадь параллелограмма на $\bar{e}_1$,
    $\bar{e}_2$ обратилась в нуль, векторы коллинеарны.

  \item \textbf{Момент удивления: вырождение базиса.} Доведите
    $\bar{e}_1 \parallel \bar{e}_2$ --- вспыхивает плашка «базис
    вырожден --- линейно зависим!», координаты $\alpha_1$, $\alpha_2$
    исчезают. Дайте паузу: только что свободно раскладывавшийся вектор
    разложить стало нельзя. Свяжите это с критерием линейной
    зависимости: теперь $\exists\,(\lambda_i)\neq 0:\ \lambda_1\bar{e}_1
    + \lambda_2\bar{e}_2 = \bar 0$.

  \item \textbf{Два лица вырождения.} При коллинеарном базисе подвигайте
    $\bar{x}$: если он вне общей прямой $\bar{e}_1$, $\bar{e}_2$ ---
    разложения \emph{не существует} (точку до него не дотянуть); если
    положить $\bar{x}$ на эту прямую --- разложений становится
    \emph{бесконечно много}. Подчеркните: именно поэтому базис обязан
    быть линейно независимым.

  \item \textbf{Связь с критерием и определителем.} Верните базис в
    невырожденное состояние и сопоставьте на доске: «независимы»
    $\iff$ «неколлинеарны» $\iff \Delta\neq 0 \iff$ «есть и однозначно
    разложение». Это четыре эквивалентные формулировки одного факта из
    Лекций~15--16.

  \item \textbf{Выход на приложения.} Завершите, переименовав на словах
    $\bar{e}_1$, $\bar{e}_2$ в «оси системы координат» или «признаки
    модели»: вырождение базиса --- это потеря оси в графике и
    мультиколлинеарность признаков в ML. Так демонстрация связывает
    геометрию плоскости с матрицей перехода, рангом и устойчивостью
    вычислений.
\end{enumerate}
